\begin{textsample}{NEG Dim 2 – vem\_ed\_29 – Score 3.00 – vem\_ed\_29/t153.txt}  \label{ex:f2_neg_125}
Osvaldo Succi Junior apresenta trabalho na Faubai 2025 Osvaldo Succi Junior, coordenador de Apoio à Internacionalização do Ensino Superior do Centro Paula Souza, apresentou trabalho na Conferência Internacional Faubai 2025, evento da Associação Brasileira de Educação Internacional (Faubai) realizado entre 26 e 30 de abril em Brasília.
Intitulada “Critical Lenses, Intelligent Tools: Exploring AI in a COIL Documentary Project” (“Lentes Críticas, Ferramentas Inteligentes: Explorando IA num projeto de documentário COIL”), a apresentação de Succi Junior narrou um Projeto Colaborativo Internacional (PCI / Cesu) realizado em 2024 entre as Fatecs Barueri e Guaratinguetá e a Universidad Católica Andrés Bello (Venezuela), voltado à questão migratória dos povos originários entre Venezuela e Brasil.
A proposta desse PCI foi estudar um aspecto comum entre os dois países, que os aproximasse.
Então, foi abordada a crise dos Yanomami, povo indígena cujo território foi dividido pelas fronteiras do Brasil e da Venezuela e que passa por situação dramática do ponto de vista humanitário. continuação Como o projeto ainda não foi encerrado, a demonstração com a ferramenta de Inteligência Artificial Generativa"NotebookLM"foi realizada com base em um conjunto de edições do informativo"VEm".
Após a análise do conteúdo de 26 edições, a ferramenta produziu um podcast de 18 minutos em inglês que surpreende pela qualidade.
Succi Junior ressaltou ainda a importância de que os educadores auxiliem os alunos a desenvolverem letramento crítico, despertando a consciência de um mundo \textbf{interconectado}, estimulando-os a estabelecer ligações do pessoal ao global, capacitando-os a analisar textos a partir de perspectivas transculturais e incentivando-os a serem social e politicamente ativos em questões “glocais” e multiculturais.
Sobre o quadro teórico do Critical Virtual Exchange proposto por Mirjam Hauck (2023), Succi Junior destacou o uso de tecnologias mais acessíveis para alunos com baixo nível socioeconômico, projetos de Intercâmbio Virtual alinhados com os Objetivos de Desenvolvimento Sustentável (ODS) da ONU e extensão com empresas, ONGs e instituições de caridade para promover o desenvolvimento de habilidades transversais, aumentar a empregabilidade dos graduados e apoiar a realização dos ODSs.
Com o tema “A Caminho de Relações Igualitárias e Sustentáveis” (Towards Equitable and Sustainable Partnerships), a Conferência Internacional Faubai 2025 reuniu representantes de instituições de ensino superior do Brasil e de outros 27 países para debater estratégias de internacionalização que promovam parcerias equitativas, sustentáveis e de benefício mútuo.
Entre os principais tópicos discutidos estiveram a internacionalização no Sul Global, a inclusão e o empoderamento, a mobilidade estudantil, as tecnologias digitais aplicadas à cooperação acadêmica e a internacionalização do currículo.
Organização fundada em 1988, a Faubai reúne cerca de 200 instituições de ensino superior brasileiras.
Dentre as principais atividades, constam a promoção e a divulgação de eventos como a Conferência Internacional, realizada anualmente desde a sua fundação.

% matched lemmas: interconectar
\end{textsample}
